\question[3]
Gegeben ist die folgende Funktion:
\begin{lstlisting}
def preis(netto, mwst=19, rabatt=0):
    return netto * (100 + mwst) // 100 - rabatt
\end{lstlisting}
Gib an, was ausgegeben wird, bzw. ob ein Fehler auftritt.
\begin{lstlisting}
print(preis(100))
print(preis(100, 7))
print(preis(100, rabatt=10))
print(preis(rabatt=5, netto=200))
print(preis(50, 7, 3))
print(preis(mwst=7))
\end{lstlisting}

\begin{solutionbox}{4cm}
\begin{lstlisting}[language={}]
119     mwst und rabatt haben ihre Standardwerte
107     mwst = 7
109     mwst = 19, rabatt = 10
233     Schlüsselwortargumente, Reihenfolge egal: 238 - 5
50      53 - 3
Fehler (TypeError): für netto gibt es keinen Standardwert
\end{lstlisting}
\end{solutionbox}
